\chapter{Jom Belajar \LaTeX{}!}
\section{Belajar \LaTeX{} susahlah!}
\label{chap1}
\LaTeX{} mempunyai cerun \index{pembelajaran} yang tinggi pada mulanya, tetapi apabila sudah dipelajari, ia akan memudahkan anda untuk menyelesaikan tugasan anda.
Sekiranya anda bergiat di dalam bidang yang memerlukan penulisan persamaan (equation) contohnya, \LaTeX{} sangat membantu anda.

\section{Mengenali format \latex}
\label{kenal}
Sekiranya anda pernah mempelajari apa-apa bahasa aturcara yang berbentuk \emph{procedural}, anda akan dapati \latex mempunyai format yang hampir serupa, di mana
turutan arahan yang akan dilaksanakan adalah dari atas ke bawah. 

\latex{} mempunyai sintaks tersendiri, di mana pengguna hendaklah mengisytiharkan awalan dan akhiran dokumen. Ia adalah seperti berikut:\\


\begin{minipage}{\textwidth}
\begin{lstlisting}[frame=shadowbox,label={cth1-arahan}]
 \begin{document}
  ...di sini anda akan laksanakan arahan anda...
 \end{document}
\end{lstlisting}

\end{minipage}


\hspace{2cm}
 %\ref{cth1-arahan} 
Seperti yang anda lihat pada Contoh di atas, itu adalah sintaks permulaan bagi dokumen \latex{} anda.

\section{Bentuk penulisan}
Secara umum, ada tiga jenis kelas dokumen yang digunakan di dalam \latex{}, iaitu 

\begin{description}
 \item[article] untuk penulisan makalah, dihantar ke seminar akademik
\item[report] mirip seperti artikel
\item[book]untuk penulisan buku, terdapat sokongan indeks, penetapan isi kandungan secara automatik, kepala dokumen dan sebagainya
\end{description}

\subsection{Kelas article}
\emph{Article class} ialah salah satu kelas dokumen yang penting dan ringkas, di mana anda perlu menetapkan kategori article ini di kepala dokumen anda.
Lihat contoh di bawah:\\

\begin{lstlisting}[frame=shadowbox]
\documentclass{article}

\begin{document}

\end{document}
\end{lstlisting}

Pada setakat ini, anda hanya telah mengisytiharkan yang dokumen anda ialah sebuah artikel. Seterusnya, kita letakkan tajuk artikel dan nama pengarang
seperti di bawah:\\

\begin{lstlisting}[frame=shadowbox]
\documentclass{article}
\title{Artikel saya}
\author{Najmi}
\begin{document}
\maketitle
\end{document}
\end{lstlisting}

\bigskip
Dokumen ini akan menghasilkan output seperti berikut:\\

\begin{minipage}{\linewidth}
\centering
\fbox{\includegraphics[scale=.7]{artikel-saya.png}}
\captionof{figure}{Artikel pertama}
\label{artikel}
\end{minipage}

\section{Bagaimana \latex{} menghasilkan dokumen?}

Secara asasnya \latex{} berfungsi seperti berikut:\\

\begin{Verbatim}[frame=single]
 fail asal (namafail.tex) ---> fail yang dijanakan (namafail.dvi)
\end{Verbatim}

Dalam kes ini, fail yang kita sunting sebagai kod sumber mempunyai sambungan .tex dan menghasilkan .dvi\index{dvi}\nomenclature{DVI}{Device Independent} . 
Tetapi, untuk memudahkan pembaca membaca dokumen yang kita hasilkan, muncullah \index{PDFLaTeX}PDFLaTeX, yang menghasilkan fail berasaskan Portable Document Format(\index{PDF}\nomenclature{PDF}{Portable Document Format} PDF).
Fail PDF boleh anda baca menggunakan pembaca PDF contohnya Acrobat Reader. Jadi dalam kes ini;\\

\begin{Verbatim}[frame=single]
 fail asal (namafail.tex) --(guna PDFLaTeX)--> fail yang dijanakan (namafail.pdf)
\end{Verbatim}

\section{Pakej di dalam \latex{}}
\latex{} memerlukan pengguna memasukkan pakej secara manual sebelum ia memproses arahan dari pengguna. Contoh-contoh pakej ialah:\\

\begin{minipage}{\textwidth}

\begin{center}
\begin{tabular}{|l|l|}
\hline
 graphicx&  untuk tujuan grafik \\ \hline
 lscape& untuk tujuan mengubah bentuk paparan kertas semasa menulis dokumen\\ \hline
 alqalam&  teks al-quran di dalam lateks\\\hline
\end{tabular}
\end{center}   
\end{minipage}
\bigskip

Jadi sebagai lanjutan untuk sebelum ini, ia akan menjadi seperti berikut:\\

\begin{lstlisting}[frame=shadowbox]
\documentclass{article}
\usepackage{graphicx}
\usepackage{lscape}
\usepackage{alqalam}

\title{Artikel saya}
\author{nama pengarang}
\begin{document}
\maketitle
\end{document}
\end{lstlisting}
